%%% FIELD target-statement
For the parameter \(K_G\) as defined, pointwise numerical edge identification on the full literal promise class \(K_G\leq 0\) is promise \(\forall\mathbb R\)-complete.  Indeed, completeness already holds on the more restricted promise subproblem
\[
\mathcal R_0
=
\bigl\{(G,\Sigma,e):(G,\Sigma,e)\text{ lies in the promise domain of }\mathsf{NUMEDGE}_3\text{ and }K_G=0\bigr\}.
\]
Thus \(\mathcal R_0\) is contained in the full literal class at \(k=0\); no equality between these two classes is asserted.  Hence \(\mathsf{ParamHandle}\) cannot yield an \(f(k)\operatorname{poly}(n)\) algorithm for the class stated in the question unless \(\mathsf P=\forall\mathbb R\).  These conclusions hold for exact fibers, including singular fibers, and use no genericity assumption.
%%% FIELD target-proof
Let \(\mathcal C_0\) denote the full literal promise class obtained by imposing \(K_G\leq 0\), without imposing any additional maximum-indegree-three restriction, and let \(\mathcal R_0\) be the subproblem displayed in the statement.  By the definition of \(K_G\),
\[
K_G=\bigl\lvert\{v\in V:\lvert\operatorname{pa}_D(v)\rvert=2\}\bigr\rvert\in\mathbb N_0.
\]
Therefore
\[
K_G\leq 0\quad\Longleftrightarrow\quad K_G=0.
\]
This equivalence identifies \(\mathcal C_0\) with the full class having no two-parent vertex; it does not impose a bound on vertices having three or more parents and hence does not identify \(\mathcal C_0\) with \(\mathcal R_0\).

Apply \(\mathrm{thm:few\mbox{-}two\mbox{-}parent\mbox{-}fixed\mbox{-}parameter\mbox{-}hardness}\) at the fixed value \(k=0\).  Its first assertion states that pointwise numerical edge identification on \(K_G\leq 0\), namely on \(\mathcal C_0\), is promise \(\forall\mathbb R\)-complete.  This proves the first assertion of the present theorem.

We next prove the separate assertion about \(\mathcal R_0\).  The same established theorem states that its hardness reduction already lands in the restriction of \(\mathsf{NUMEDGE}_3\) having \(K_G=0\), rational strictly diagonally dominant positive-definite covariance, and a depth-one directed graph of maximum directed outdegree one whose directed weak components are isolated vertices, single directed edges, or three-parent in-stars.  Every output of that reduction therefore belongs to \(\mathcal R_0\).  Because the source promise problem is promise \(\forall\mathbb R\)-hard and the reduction is promise preserving by the cited theorem, \(\mathcal R_0\) is promise \(\forall\mathbb R\)-hard.

For membership, the completeness assertion for \(\mathcal C_0\) supplies a polynomial-size universal real predicate for its YES language.  Write that predicate as
\[
U(x)\equiv \forall y\in\mathbb R^r\;\Phi(x,y),
\]
where \(x\) is the finite binary encoding fixed by \(\mathrm{def:promise\mbox{-}problem}\), \(r\) is polynomially bounded in \(\lvert x\rvert\), and \(\Phi\) is the corresponding polynomial-time constructible quantifier-free formula over rational polynomial equalities and inequalities.  Let \(b(x)\) be the Boolean condition that the encoded directed graph has maximum directed indegree at most three and that \(K_G=0\).  Both parts of \(b(x)\) are decidable in polynomial time by counting the incoming entries in each row of the encoded directed adjacency matrix.  Since \(b(x)\) contains no real variable,
\[
b(x)\wedge U(x)
\quad\Longleftrightarrow\quad
\forall y\in\mathbb R^r\;\bigl(b(x)\wedge\Phi(x,y)\bigr).
\]
Concretely, the polynomial-time formula constructor first evaluates \(b(x)\): if it is true, it outputs \(U(x)\), and if it is false, it outputs the fixed false universal sentence \(\forall y\;(0=1)\).  This realizes the displayed conjunction using only a quantifier-free rational polynomial formula and has polynomial encoding length.  On the promise domain of \(\mathcal R_0\), \(b(x)\) holds and \(U(x)\) is exactly the numerical-identification predicate; outside that promise domain its value is immaterial.  Thus this construction is a valid promise \(\forall\mathbb R\) representation of \(\mathcal R_0\).  Combining membership with the preceding hardness proves that \(\mathcal R_0\) is promise \(\forall\mathbb R\)-complete.  Notice that this argument proves containment and completeness separately and nowhere replaces the full class \(\mathcal C_0\) by \(\mathcal R_0\).

For the algorithmic consequence, suppose that \(\mathsf{ParamHandle}\) yielded an exact algorithm for the full literal parameterized class with bit complexity \(f(k)\operatorname{poly}(n)\).  On \(\mathcal C_0\) one has \(k=0\), so its running time is \(f(0)\operatorname{poly}(n)\).  If \(N\) is the total binary input length, then \(n\leq N\), and hence this is polynomial in \(N\).  It would decide the promise \(\forall\mathbb R\)-complete problem on \(\mathcal C_0\) in polynomial time.  Equivalently, this is the \(k=0\) case of the complexity consequence already established by \(\mathrm{thm:few\mbox{-}two\mbox{-}parent\mbox{-}fixed\mbox{-}parameter\mbox{-}hardness}\).  Therefore such an algorithm would imply \(\mathsf P=\forall\mathbb R\).

Finally, \(\mathrm{thm:few\mbox{-}two\mbox{-}parent\mbox{-}fixed\mbox{-}parameter\mbox{-}hardness}\) states that its conclusions concern exact coefficient fibers, including singular fibers, and require no genericity assumption.  Specializing \(k\) to zero and restricting its explicit hard family to \(\mathcal R_0\) changes neither fact.  This proves the final assertion.
